\begin{corollary}[K17-K12-R8-INDEX5]
\label{cor:k12-r8-index-five}
Every potentially contributing rank-eight projective degree-twelve
Voronoi cell has exactly thirteen minimum pairs, and its maximal
minimum-basis index is at most five. More precisely, for each minimum
frame $F=\sum_{i=1}^8\mathbb Zf_i$, the quotient $L/F$, up to signs,
permutations of the frame and changes of quotient generators, is one
of the following types. Superscripts on entries denote repetition.
\[
\begin{array}{c|l}
 L/F & \text{coefficients of generators on the frame}\\ \hline
 0 & \text{none}\\
 C_2 &(1^m,0^{8-m})/2,\quad 5\leq m\leq8\\
 C_3 &(1^m,0^{8-m})/3,\quad 7\leq m\leq8\\
 C_4 &(1^5,2^2,0)/4,\ (1^5,2^3)/4,\
       (1^6,2^2)/4,\ (1^7,2)/4\\
 C_2\times C_2 &(1,1,1,1,1,0,0,0)/2,\
                 (0,0,0,1,1,1,1,1)/2\\
 C_5 &(1^5,2^3)/5
\end{array}
\]
This is a list of necessary quotient types, including possible
ambient glue, rather than a classification of the complete cells.
\end{corollary}

\begin{proof}
Theorem~\ref{thm:k12-r8-nonsimplicial} removes the nonsimplicial
cells. A remaining cell has thirteen independent quadratic rays,
hence exactly thirteen minimum pairs.

Use the complete minimum-frame quotient classification in
\cite[\S11 and Table~11.1, author PDF pp.~45--48]{MartinetIndex2001}.
Its setup specifies that each type has a unique minimal minimum
configuration and that the column $s$ records its number of pairs.
A type supported on $n\leq8$ frame coordinates forces at least
\[
                           s+(8-n)                     \tag{R8.6}
\]
minimum pairs in the ambient rank-eight lattice: the support
lattice already contains $s$ pairs and the remaining frame vectors
supply $8-n$ further pairs.

The completed index bound, from
\cite[appendix]{MartinetIndexErrata2014} and
\cite[Theorem~10.1]{MartinetIndex2001}, excludes orders above eight
outside $E_8$; that exception has $120$ pairs. The complete quotient
list contains no order seven. The six index-six types first occur
in rank eight and have pair counts
\[
                     31,\ 27,\ 120,\ 28,\ 14,\ 36.
\]
For index eight, the primitive rank-seven type is $E_7$ with
$63$ pairs, and the six primitive rank-eight types have counts
\[
                     48,\ 75,\ 120,\ 56,\ 48,\ 32.
\]
All exceed the allowable count in (R8.6). Thus every minimum-basis
index is at most five.

For the remaining orders, apply (R8.6) to the same complete table.
Order two with support four forces $12+4=16$ pairs, leaving
supports five through eight. Order three with support six forces
$12+2=14$, leaving supports seven and eight. At order four, every
row except the four cyclic types and one elementary type displayed
above forces at least sixteen pairs. The first displayed cyclic
type has support seven, and the other four types have support
eight. At order five, the distributions $(4,4)$ and $(6,2)$ force
sixteen pairs, leaving $(5,3)$. This gives precisely the stated
necessary list.

For the signed-coordinate assertion in the order-five row, one can
use the direct calculation in \cite[\S9, case $d=5$, p.~31;
Proposition~9.1 and proof, p.~33]{MartinetIndex2001}.
After signs and a change of the single quotient generator, it lists
exactly the distributions $(6,2),(5,3),(4,4)$ and proves that the
first and last force sixteen pairs. Thus this assertion does not
rely on interpreting the broader monomial-code equivalence in
Definition~7.4 as a signed change of frame.
\end{proof}

No assertion is made that each remaining quotient type extends to
a cell with exactly thirteen rays, or that the remaining cells or
their homology classes vanish. The published quotient-type
classification is an explicit retained input.
